\documentclass[12pt,a4paper]{article}

\usepackage[utf8]{inputenc}
\usepackage[T1]{fontenc}
\usepackage[margin=2.5cm]{geometry}
\usepackage{hyperref}
\usepackage{microtype}
\usepackage{parskip}
\usepackage{lmodern}

\hypersetup{colorlinks=true, urlcolor=blue!60!black}

\pagestyle{empty}

\begin{document}

\begin{flushright}
Kundai Farai Sachikonye \\
Zeltnerstrasse 30, 90443 Nürnberg, Germany \\
\href{https://github.com/fullscreen-triangle}{github.com/fullscreen-triangle} \\
\texttt{kundai.sachikonye@bitspark.com} \\
(+49) 1626386344 \\[6pt]
\today
\end{flushright}

\vspace{1em}

Prof. Dr. Peter Robin Hiesinger \\
Freie Universität Berlin \\
Fachbereich Biologie, Chemie, Pharmazie --- Institut für Biologie \\
Königin-Luise-Str. 1--3 \\
14195 Berlin (Dahlem)

\vspace{1.5em}

\textbf{Re: Wiss.\ Mitarbeiter*in, AG Neurobiologie --- Kennung 21259000\_2026}

\vspace{1em}

Dear Professor Hiesinger,

I am writing regarding the research-associate position in your group. The
position explicitly includes the group's complementary computational work ---
digital image analysis, statistical evaluation of imaging data, and transcriptome
and connectome analysis --- and that is precisely where my work sits. My
background is computational, and it is on that side that I would contribute to
your group's imaging and analysis of Drosophila brains.

\textbf{Digital image analysis of microscopy data.}
This is my strongest point of overlap. I have built a quantitative
image-analysis framework for fluorescence and phase-contrast microscopy ---
segmentation, point-spread-function deconvolution, scale-field and coordinate
reconstruction, and morphometry with uncertainty quantification, GPU-accelerated
and available as a live tool
(\href{https://hieronymus-omega.vercel.app/}{hieronymus-omega.vercel.app}). For a
group whose technical focus is live imaging of Drosophila brains, rigorous,
reproducible image quantification is something I could contribute immediately.

\textbf{Connectome and computational neuroscience.}
My independent research includes neural-circuit modelling: a closed-circuit
decomposition of motor and neural signals into supraspinal and peripheral
contributions, applied to neurodegeneration (Parkinson's, ataxia), together with
circuit-graph (connectome-style) network analysis. This is directly relevant to
the group's connectome and computational-modelling work.

\textbf{Transcriptome analysis.}
I have developed pipelines integrating transcriptomics with genotype and
metagenomics, validated against 1000~Genomes, gnomAD, and UK~Biobank.

\textbf{How I would contribute.}
My strength is the quantitative and computational side: turning imaging and
sequencing data into rigorous, reproducible, statistically sound results, and
building the tools to do so. I work closely and naturally with experimental
groups, complementing their imaging and genetics rather than duplicating it.
English is native; I have been resident in Germany since 2011.

If the computational and image-analysis side of the position is one you would
consider filling with a computational specialist, I would welcome a conversation.

\vspace{1.5em}
Yours sincerely,

\vspace{2.5em}
\textbf{Kundai Farai Sachikonye}

\end{document}
